% C5 — The diagonal Amos-type family: quantitative crossing
% classification, machine-checked.  Charter: docs/C5-CHARTER.md
% (Amendments 1-8).  REGIME: tricotomy; related work per the
% Amendment-7 directed collation VERBATIM substance; novelty ONLY
% on the crossing package; RS16 prior art declared expressly;
% same-lane ids pending per the C4 release precedent.
\documentclass[11pt]{article}
\usepackage[margin=1.1in]{geometry}
\usepackage{amsmath,amssymb,amsthm}
\usepackage{booktabs}
\usepackage[colorlinks=true,linkcolor=blue,urlcolor=blue,citecolor=blue]{hyperref}

\newtheorem{theorem}{Theorem}[section]
\newtheorem{lemma}[theorem]{Lemma}
\newtheorem{remark}[theorem]{Remark}

% long monospace artifact names offer no hyphenation points;
% let paragraphs stretch rather than overflow the margin
\emergencystretch=3em
\tolerance=2000

\newcommand{\lean}[1]{\texttt{#1}}
\newcommand{\Bfam}{\mathrm{B}}

\title{The diagonal Amos-type family at real order:\\
a machine-checked quantitative crossing classification}
\author{Lluis Eriksson\thanks{This manuscript and parts of the formal
development were produced with AI assistance under the repository's
audit protocol.  Every theorem below is machine-checked in Lean~4
against a pinned Mathlib; the independent numerical companion of
Section~\ref{sec:companion} is a committed interval-arithmetic
transcript; remaining numerics are labelled paper-level where they
occur.  Repository:
\url{https://github.com/lluiseriksson/THE-ERIKSSON-PROGRAMME}.}}
\date{July 2026}

\begin{document}
\maketitle

\begin{abstract}
For the one-parameter family of Amos-type expressions
\[
\Bfam_{\nu,c}(x) \;=\; \frac{x}{\nu+c+\sqrt{(\nu+c)^2+x^2}},
\qquad c \in \mathbb{R},
\]
whose member $c = \tfrac12$ is the classical Amos-type upper bound
for the modified Bessel ratio $\rho_\nu = I_{\nu+1}/I_\nu$, we
formalize in Lean~4, at every real order $\nu \ge 0$ over the
$\Gamma$-power series, the classification of the parameter:
$\Bfam_{\nu,c}$ is a uniform upper bound for $\rho_\nu$ (all
$\nu \ge 0$, $x > 0$) exactly when $c \le \tfrac12$, and a uniform
lower bound for every $c \ge 1$, with explicit counterexample
witnesses.  The classification statement itself is
equivalent to known results of Ruiz-Antol\'in and Segura; for it we
claim only the machine-checking.  The contribution is the regime
between the ends: for every $\nu \ge 0$ and every
$c \in (\tfrac12, 1)$ we prove that the fixed diagonal member
$\Bfam_{\nu,c}$ crosses $\rho_\nu$ \emph{exactly once} on
$(0,\infty)$ --- a transversal crossing located in the explicit
window $[\sqrt{1-c},\, 2\hat{D}/(2c-1)]$ with
$\hat{D} = (\nu+\tfrac32)^2 - (\nu+c)^2$, obeying the strict
threshold $x_* > x_\dagger = 2(\nu+c)\sqrt{c(1-c)}/(2c-1)$, with
globally determined sign on both sides and the two-sided scale law
$2(\nu+c)\sqrt{c(1-c)} \le (2c-1)\,x_* \le 2\hat{D}$.  Degenerate
(critical) contact is \emph{excluded by theorem}, through the exact
identity $(\nu+c)^2\,D''(x_\dagger) = (2c-1)^3\,\Bfam_{\nu,c}(x_\dagger)$
for the gap $D = \rho_\nu - \Bfam_{\nu,c}$ at any hypothetical
critical contact.  The entire chain carries the
axiom oracle
$[\lean{propext}, \lean{Classical.choice}, \lean{Quot.sound}]$ and
no analytic hypothesis beyond $\nu \ge 0$, $x > 0$; a pre-registered
certified interval-arithmetic companion verifies the crossing
phenomenon independently of the C5 crossing theorems at $30$ parameter
pairs spanning the hard regimes, all passing at $128$ bits.
\end{abstract}

\section{Introduction}

The ratio $\rho_\nu(x) = I_{\nu+1}(x)/I_\nu(x)$ of modified Bessel
functions of the first kind admits algebraically calibrated bounds
of the shape $x/(\lambda+\sqrt{\mu^2+x^2})$; the literature on such
bounds --- Amos \cite{Amos74}, Segura
\cite{Segura11,Segura21,Segura23}, Hornik--Gr\"un \cite{HG13,HG24},
Ruiz-Antol\'in--Segura \cite{RS16} --- is rich, and the
one-parameter \emph{diagonal} family studied here,
\[
\Bfam_{\nu,c}(x) \;=\; \frac{x}{\nu+c+\sqrt{(\nu+c)^2+x^2}},
\]
is, up to reparametrization, the family whose sharp ends
\cite{RS16} identifies: under $\mu = \nu+1$, $\alpha = 2c-1$, the
family $b_\alpha$ of \cite{RS16} is exactly $\Bfam_{\nu,c}$.
\emph{No novelty is claimed for the family, nor for the qualitative
observation that its non-bounding members must cross the ratio}
--- \cite{RS16} already notes that for parameters between the sharp
ends the graphs cross.  Three companion papers of this series
formalized the consequence calculus of the $c=\tfrac12$ bound at
integer order \cite{ErikssonC2}, the bound itself at integer order
\cite{ErikssonC3}, and the bound at every real order $\nu \ge 0$
over the $\Gamma$-power series \cite{ErikssonC4}; the present paper
builds directly on the last.

The contribution here is twofold, and the two halves carry
different claims:
\begin{enumerate}
\item \textbf{The classification, machine-checked.}  For every real
$\nu \ge 0$ and $x > 0$: $\rho_\nu < \Bfam_{\nu,c}$ holds uniformly
iff $c \le \tfrac12$ (Theorem~\ref{thm:classification}), and for
$c \ge 1$ the reverse inequality holds uniformly
(Theorem~\ref{thm:lowerhalf}).  As mathematics this is equivalent
to results of \cite{RS16}; the claim is the formalization ---
uniform in real order, from the series, with explicit rational
witnesses, inside one pinned axiom-clean development.
\item \textbf{The quantitative crossing theorem.}  For $\nu \ge 0$
and $\tfrac12 < c < 1$, the fixed diagonal member $\Bfam_{\nu,c}$
crosses $\rho_\nu$ exactly once; the crossing is transversal,
satisfies an explicit strict threshold, determines the sign of
$\rho_\nu - \Bfam_{\nu,c}$ globally on both sides, lies in a
constructive finite window, and obeys an explicit two-sided scale
law (Theorem~\ref{thm:crossing}).  Degenerate contact is excluded
by an exact second-derivative identity
(Theorem~\ref{thm:nocontact}).  We are not aware of a previous
theorem giving this complete quantitative crossing classification
for the diagonal family; Section~\ref{sec:related} compares against
the closest antecedents parameter by parameter.
\end{enumerate}

Statuses: \textbf{exact} means a Lean theorem in the pinned central
repository with clean axiom oracle; \textbf{certified} means a
committed interval-arithmetic transcript; \textbf{paper-level}
means ordinary prose, always flagged; \textbf{verified} marks an
ordinary high-precision computation used only as a cross-check of a
certified result, never as evidence.  The repository's largest
namespace (\lean{YangMills/}) reflects its original motivation in
lattice gauge theory; no reduction to any continuum or
gauge-theoretic problem is claimed or used.

\section{Setting}\label{sec:setting}

All objects live in module \lean{AmosClosure} of the pinned central
repository, over the real-order interface of \cite{ErikssonC4}:
$I_\nu = \lean{besselIReal}\ \nu$ is the $\Gamma$-power series
$\sum_{k\ge0}(x/2)^{\nu+2k}/(k!\,\Gamma(\nu+k+1))$ (real exponents
via \lean{rpow}), with positivity, the three-term recurrence,
termwise differentiability, the Riccati equation
\begin{equation}\label{eq:riccati}
\rho_\nu'(x) \;=\; Q_x(\rho_\nu(x)), \qquad
Q_x(y) := 1 - \frac{2\nu+1}{x}\,y - y^2,
\end{equation}
the barrier $\psi_\nu > 2\nu+1$ for
$\psi_\nu(x) = x(1/\rho_\nu - \rho_\nu)$, and the $c=\tfrac12$
bound $\rho_\nu < \Bfam_{\nu,1/2}$ (\lean{amosBoundReal\_holds}).
New objects of this paper: the family
$\Bfam_{\nu,c} = \lean{amosFamily}\ \nu\ c$, the lower bound
$L_\nu = \lean{amosLower}\ \nu$ of Section~\ref{sec:lower}, and the
crossing gap
\[
D(x) \;=\; \rho_\nu(x) - \Bfam_{\nu,c}(x).
\]
Throughout, $a = \nu+c$ and $s = \sqrt{a^2+x^2}$.  Everything is
stated at $x > 0$, the \lean{rpow} domain of the development.

\section{Main results}\label{sec:results}

\begin{theorem}[classification of the upper-bound range;
\lean{amosFamily\_uniform\_upper\_iff}]\label{thm:classification}
For every $c \in \mathbb{R}$:
\[
\Bigl(\forall\, \nu \ge 0,\ \forall\, x > 0:\;
\rho_\nu(x) < \Bfam_{\nu,c}(x)\Bigr)
\;\Longleftrightarrow\;
c \le \tfrac12 .
\]
The failure witness for $c > \tfrac12$ is explicit and rational:
$x_0(c) = 2/(2c-1)$ at $\nu = 0$
(\lean{amosFamily\_lt\_amosLower\_witness}).
\end{theorem}

\begin{theorem}[the uniform lower half;
\lean{amosFamily\_lower\_of\_one\_le}]\label{thm:lowerhalf}
For every $c \ge 1$, every $\nu \ge 0$ and every $x > 0$,
$\Bfam_{\nu,c}(x) < \rho_\nu(x)$.
\end{theorem}

Theorems \ref{thm:classification} and \ref{thm:lowerhalf} are
machine-checked forms of known mathematics \cite{RS16}; between
them lies the regime $\tfrac12 < c < 1$, where $\Bfam_{\nu,c}$
bounds on neither side.  There --- with the strictness of the
threshold in item~(ii), and through it the unconditional uniqueness
and orientation, delivered by Theorem~\ref{thm:nocontact} below:

\begin{theorem}[the global crossing theorem;
\lean{amosFamily\_global\_crossing},
\lean{amosFamily\_crossing\_orientation}]\label{thm:crossing}
Fix $\nu \ge 0$ and $\tfrac12 < c < 1$.  There is exactly one
$x_* > 0$ with $\rho_\nu(x_*) = \Bfam_{\nu,c}(x_*)$, and:
\begin{enumerate}
\item[(i)] \emph{(window)} $\sqrt{1-c} \;\le\; x_* \;\le\;
\dfrac{2\hat{D}}{2c-1}$, where
$\hat{D} = (\nu+\tfrac32)^2-(\nu+c)^2$
(\lean{amosFamily\_crossing\_exists});
\item[(ii)] \emph{(strict threshold)} $x_* > x_\dagger :=
\dfrac{2(\nu+c)\sqrt{c(1-c)}}{2c-1}$
(\lean{crossing\_threshold\_strict});
\item[(iii)] \emph{(transversality)} at the crossing,
$D'(x_*) = \dfrac{\Bfam_{\nu,c}(x_*)}{x_*}
\Bigl(2c-1-\dfrac{\nu+c}{\sqrt{(\nu+c)^2+x_*^2}}\Bigr) > 0$
(\lean{crossing\_hasDerivAt}, whose factor is positive by
\lean{crossing\_threshold\_strict});
\item[(iv)] \emph{(orientation)} $\rho_\nu < \Bfam_{\nu,c}$ on
$(0, x_*)$ and $\rho_\nu > \Bfam_{\nu,c}$ on $(x_*, \infty)$
(\lean{amosFamily\_crossing\_orientation}); conversely --- an
immediate trichotomy consequence --- the sign of $D$ locates $x$
relative to $x_*$;
\item[(v)] \emph{(scale law)} $2(\nu+c)\sqrt{c(1-c)} \;\le\;
(2c-1)\,x_* \;\le\; 2\hat{D}$ (\lean{crossing\_scale}).
\end{enumerate}
\end{theorem}

\begin{theorem}[no degenerate contact;
\lean{crossing\_no\_critical\_contact}]\label{thm:nocontact}
Fix $\nu \ge 0$ and $\tfrac12 < c < 1$.  No crossing of $\rho_\nu$
and $\Bfam_{\nu,c}$ is a critical point of $D$: at any hypothetical
contact with $D = D' = 0$ --- which the slope identity forces to
lie exactly at $x_\dagger$ --- the second derivative obeys the
exact identity
\[
(\nu+c)^2\; D''(x_\dagger) \;=\; (2c-1)^3\;
\Bfam_{\nu,c}(x_\dagger) \;>\; 0,
\]
and a critical contact with positive curvature contradicts the
below-threshold orientation.  Consequently the threshold in
Theorem~\ref{thm:crossing} is strict and uniqueness and orientation
hold unconditionally.
\end{theorem}

\begin{theorem}[the Tur\'an inequality at real order;
\lean{besselIReal\_turan}]\label{thm:turan}
For every $\nu \ge 0$ and $x > 0$,
$I_{\nu+1}(x)^2 > I_\nu(x)\, I_{\nu+2}(x)$; equivalently
$\psi_\nu(x) < 2\nu+2$, completing the sandwich
$2\nu+1 < \psi_\nu(x) < 2\nu+2$ whose lower half is the barrier of
\cite{ErikssonC4}.  This inequality is classical (Tur\'an-type
inequalities for $I_\nu$ go back to Thiruvenkatachar--Nanjundiah;
see Baricz--Ponnusamy \cite{BP13}); it is proved here at real order
over the $\Gamma$-series and used as supporting structure.
\end{theorem}

\section{The proof}\label{sec:proof}

The formal proofs are the authority; this section presents the
mathematics in full, at the level of the checklist a referee needs.

\subsection{Family algebra: calibration and the nullcline
trichotomy}\label{sec:family}

Since $s = \sqrt{a^2+x^2} > |a| \ge -a$, the denominator $a+s$ is
positive for \emph{every} real $c$, and the conjugate form
$\Bfam_{\nu,c} = (s-a)/x$ follows from $(s-a)(s+a) = x^2$
(\lean{amosFamily\_pos}).  The general calibration identity
\begin{equation}\label{eq:calibration}
\frac{1}{\Bfam_{\nu,c}(x)} - \Bfam_{\nu,c}(x) \;=\;
\frac{2(\nu+c)}{x}
\qquad (\lean{amosFamily\_calibration})
\end{equation}
extends the $c = \tfrac12$ calibration of \cite{ErikssonC4}.
Substituting into the Riccati quadratic of \eqref{eq:riccati}:
$1 - \Bfam^2 = (2a/x)\Bfam$, so
\begin{equation}\label{eq:nullcline}
Q_x\bigl(\Bfam_{\nu,c}(x)\bigr) \;=\;
\frac{2a - (2\nu+1)}{x}\,\Bfam_{\nu,c}(x) \;=\;
\frac{2c-1}{x}\,\Bfam_{\nu,c}(x)
\qquad (\lean{riccatiQReal\_amosFamily}),
\end{equation}
with the sign trichotomy: zero iff $c = \tfrac12$ (the member
\emph{on} the nullcline --- exactly the calibration that drives the
proof of the $c=\tfrac12$ bound), positive iff $c > \tfrac12$,
negative iff $c < \tfrac12$.  Since
$a \mapsto x/(a+\sqrt{a^2+x^2})$ is decreasing
(\lean{amosFamily\_anti}; the non-strict form is all the chain
uses), $c \le \tfrac12$ gives
$\Bfam_{\nu,c} \ge \Bfam_{\nu,1/2} > \rho_\nu$ by
\cite{ErikssonC4}: the sufficiency half of
Theorem~\ref{thm:classification}
(\lean{amosFamily\_upper\_of\_le\_half}).

\subsection{The lower bound by recurrence inversion}\label{sec:lower}

The ratio recurrence
$1/\rho_\nu = 2(\nu+1)/x + \rho_{\nu+1}$
(\lean{besselIReal\_ratio\_recurrence}) turns the $c=\tfrac12$
\emph{upper} bound at order $\nu+1$ into a \emph{lower} bound at
order $\nu$: with $b = \nu+\tfrac32$ and $s' = \sqrt{b^2+x^2}$,
\[
\frac{1}{\rho_\nu(x)} \;<\; \frac{2(\nu+1)}{x} + \frac{s'-b}{x}
\;=\; \frac{\nu+\tfrac12+s'}{x},
\quad\text{i.e.}\quad
\rho_\nu(x) \;>\; \frac{x}{\nu+\tfrac12+\sqrt{(\nu+\tfrac32)^2+x^2}}
\;=:\; L_\nu(x),
\]
where the middle step is the algebraic collapse
$x^2/(b+s') = s'-b$ (\lean{besselLowerReal\_holds}).  This is the
sharp-at-infinity lower bound of \cite{RS16}; here it is a
three-line corollary of the real-order upper bound, which is the
point of recording it.  $L_\nu$ matches $\rho_\nu$ to first order
at infinity ($L_\nu \sim 1 - (\nu+\tfrac12)/x$), while
$\Bfam_{\nu,c} \sim 1 - (\nu+c)/x$ (paper-level asymptotics, not
used in the proofs); for $c > \tfrac12$ the family
member falls below $L_\nu$ at explicit finite $x$: at $\nu = 0$ and
$x_0(c) = 2/(2c-1)$ the inequality
$\Bfam_{0,c}(x_0) < L_0(x_0)$ reduces, after clearing the
calibrated denominators with $(2c-1)x_0 = 2$, to the rational
inequality $\tfrac94 - c^2 < 2$ for $c \in (\tfrac12, 1]$ and its
routine extension beyond (\lean{amosFamily\_lt\_amosLower\_witness});
chaining $\Bfam_{0,c}(x_0) < L_0(x_0) < \rho_0(x_0)$ kills
uniformity.  That is the necessity half of
Theorem~\ref{thm:classification}.  For Theorem~\ref{thm:lowerhalf}:
$c \ge 1$ gives $\Bfam_{\nu,c} \le \Bfam_{\nu,1} < L_\nu < \rho_\nu$,
where $\Bfam_{\nu,1} < L_\nu$ is the root comparison
$\sqrt{(\nu+\tfrac32)^2+x^2} < \tfrac12 + \sqrt{(\nu+1)^2+x^2}$,
i.e.\ $\nu+1 < \sqrt{(\nu+1)^2+x^2}$ after squaring
(\lean{amosFamily\_one\_lt\_amosLower}).

\subsection{The Tur\'an sandwich}\label{sec:turan}

$\rho_\nu > \Bfam_{\nu,1}$ (from Section~\ref{sec:lower}) plus the
calibration \eqref{eq:calibration} at $c = 1$ gives
$\psi_\nu < 2(\nu+1)$ (\lean{besselPsiReal\_lt}); together with the
barrier $\psi_\nu > 2\nu+1$ of \cite{ErikssonC4} this is the
sandwich $2\nu+1 < \psi_\nu < 2\nu+2$.  The upper half converts,
through the recurrence, into $\rho_\nu > \rho_{\nu+1}$ and thence
into the Tur\'an inequality of Theorem~\ref{thm:turan}
(\lean{besselIReal\_turan}).  We emphasize provenance: the
mathematical content of the sandwich upper half is classical
\cite{BP13}; the formal statement at real order over the
$\Gamma$-series is what is new here, and the crossing results below
do not depend on it.

\subsection{Existence: the explicit window}\label{sec:window}

Fix $\tfrac12 < c < 1$.  At the left end, the product-form upper
bound $\rho_\nu(x) < x/\bigl(2(\nu+1)(1-q_\nu)\bigr)$ with
$q_\nu = (x/2)^2/(\nu+2)$, valid for $q_\nu < 1$
(\lean{besselRatioReal\_lt\_product},
from the geometric-domination toolkit of \cite{ErikssonC4}) beats
the family member below $\sqrt{1-c}$: using
$\sqrt{a^2+x^2} \le a + x^2/(2a)$ one checks
$D(x) < 0$ for all $0 < x \le \sqrt{1-c}$, uniformly in $\nu$
(\lean{crossing\_witness\_small\_le}).  At the right end, at
$X = 2\hat{D}/(2c-1)$ with $\hat{D} = (\nu+\tfrac32)^2-(\nu+c)^2$,
the two calibrated roots separate by less than $(2c-1)/4$ and
$L_\nu(X) > \Bfam_{\nu,c}(X)$, so $D(X) > 0$
(\lean{crossing\_witness\_large}).  Continuity of $D$ (from
differentiability) and the intermediate value theorem give a
crossing in the window (\lean{amosFamily\_crossing\_exists}).

\subsection{Slope identity, threshold, uniqueness,
orientation}\label{sec:slope}

The family derivative is
$\Bfam_{\nu,c}'(x) = a/\bigl(s\,(a+s)\bigr)$
(\lean{amosFamily\_hasDerivAt}).  At any zero of $D$, $\rho_\nu$
takes the family value, so by \eqref{eq:riccati} and
\eqref{eq:nullcline} $\rho_\nu' = \tfrac{2c-1}{x}\Bfam_{\nu,c}$
there, while $a/(s(a+s)) = a\Bfam_{\nu,c}/(sx)$ by the conjugate
form; hence the \emph{slope identity}
\begin{equation}\label{eq:slope}
D'(\text{zero}) \;=\;
\frac{\Bfam_{\nu,c}}{x}\Bigl(2c-1-\frac{a}{s}\Bigr)
\qquad (\lean{crossing\_hasDerivAt}).
\end{equation}
The factor $2c-1-a/s$ is negative below, zero at, and positive
above $x_\dagger = 2a\sqrt{c(1-c)}/(2c-1)$ (where
$s = a/(2c-1)$).  Three consequences, in order: every crossing
satisfies $x \ge x_\dagger$, since a crossing strictly below the
threshold would be a strictly \emph{downward} crossing yet $D < 0$
must persist from the left window endpoint
(\lean{crossing\_threshold}); above the threshold every crossing is
strictly upward (\lean{crossing\_pos\_right}), and two strictly
upward zeros of a differentiable function cannot occur without an
intermediate downward zero --- excluded --- so the crossing is
unique (\lean{crossing\_unique}); and the sign of $D$ is then
globally determined on both sides
(\lean{crossing\_orientation\_after},
\lean{crossing\_orientation\_below}).

\subsection{Exclusion of the degenerate contact}\label{sec:tangency}

One case survives Section~\ref{sec:slope}: a contact exactly at
$x_\dagger$, where the slope factor vanishes identically and the
crossing would be a critical point of $D$.  At such a contact every
quantity is an explicit algebraic function of $(\nu, c)$: writing
$t = 2c-1 \in (0,1)$,
\[
s = \frac{a}{t}, \qquad
\rho_\nu = \Bfam_{\nu,c} = \sqrt{\frac{1-c}{c}}, \qquad
\Bfam_{\nu,c}\, x = s-a = \frac{a(1-t)}{t}, \qquad
x^2 = \frac{a^2(1-t)(1+t)}{t^2},
\]
and $2\nu+1 = 2a-t$.  Differentiating the Riccati equation along
the flow and the family derivative once more
(\lean{crossing\_hasDerivAt\_everywhere} provides $D$'s derivative
as an explicit function everywhere; that function is differentiated
once more inside the tangency proof),
the two curvatures collapse to
\[
\rho_\nu'' \;=\; \frac{\Bfam}{x^2}\bigl[(2\nu+1)(1-t) -
2t\,\Bfam x\bigr] \;=\; -\frac{t^3\,\Bfam}{a^2(1+t)},
\qquad
\Bfam'' \;=\; -\frac{a\,(x/s)(a+2s)}{s^2(a+s)^2}
\;=\; -\frac{t^3\,\Bfam\,(t+2)}{a^2(t+1)},
\]
whence the exact identity of Theorem~\ref{thm:nocontact}:
$D'' = t^3\Bfam/a^2$, i.e.
$(\nu+c)^2 D'' = (2c-1)^3\,\Bfam_{\nu,c} > 0$.  In the formal
development the identity is proved division-free in exactly this
form; positivity needs only $2c-1 > 0$, $\nu+c > 0$,
$\Bfam_{\nu,c} > 0$.  The contradiction: at a critical contact,
$D' = 0$ and $D'' > 0$ force $D' < 0$ on a left interval, so $D$ is
strictly decreasing into the contact and $D(y) > D(x_\dagger) = 0$
just left of it; but $y < x_\dagger$ is strictly below the
threshold, where the orientation of Section~\ref{sec:slope} gives
$D(y) < 0$ (\lean{crossing\_no\_critical\_contact}).  No critical
contact exists; the threshold is strict
(\lean{crossing\_threshold\_strict}); uniqueness and orientation
are unconditional (\lean{crossing\_unique\_unconditional}); and the
packaged existence-and-uniqueness statement is
\lean{amosFamily\_global\_crossing}.  The scale law of
Theorem~\ref{thm:crossing}(v) is then immediate from the threshold
and the window (\lean{crossing\_scale}).

\section{Scope}\label{sec:scope}

All results are for the in-core $\Gamma$-power-series definition of
$I_\nu$ at real order $\nu \ge 0$, $x > 0$; no identification with
an external special-functions library object is claimed (none
exists at the pinned Mathlib to identify with), though the
identification theorem of \cite{ErikssonC4} ties the object to the
factorial series at every integer order.
Theorems~\ref{thm:classification} and \ref{thm:lowerhalf} are
machine-checked forms of known mathematics \cite{RS16}, and the
crossing \emph{phenomenon} for intermediate parameters is likewise
observed there; the claims of this paper are (i) the
formalization, uniform in real order, and (ii) the quantitative
crossing package of Theorems~\ref{thm:crossing} and
\ref{thm:nocontact}: unique transversal crossing, strict explicit
threshold, bilateral orientation, constructive window, two-sided
scale law, degenerate contact excluded.  We are not aware of a
previous theorem giving this complete quantitative crossing
classification for the diagonal family.

\section{The certified companion}\label{sec:companion}

A certified interval-arithmetic companion verifies the crossing
phenomenon \emph{independently of the C5 crossing theorems} ---
without consulting any C5 crossing, uniqueness, orientation, or
scale-law theorem: its protocol
--- registered in the development's charter in its own commits,
strictly before the certified run and before any certified result
was observed --- mandates that verdicts be decided \emph{only} by ball
signs of $D$ and $D'$, never by consulting the crossing theorems;
that the bisection search interval come from the fixed rule
$[\,(1-c)/4,\; 8(\nu+2)^2/(2c-1)\,]$, deliberately \emph{wider}
than the Lean window, with the theoretical window checked
afterwards rather than searched; and that failures escalate a
precision ladder ($128$ to $4096$ bits) rather than widen
intervals.  The grid of $30$ pairs,
$\nu \in \{0, \tfrac12, 1, \pi, 10, 100\} \times
c \in \{0.501, 0.6, 0.75, 0.9, 0.999\}$ ($\pi$ a certified ball),
forces the hard regimes: $c \downarrow \tfrac12$ (the crossing
escapes to infinity --- at $(\nu,c) = (100, 0.501)$ the search
interval reaches $4.2\times10^7$), $c \uparrow 1$ (the crossing
value $\Bfam(x_\dagger) \to 0$), $\nu = 0$, large order, and the
doubly-extreme corners.  Per pair, four certificates, all by ball
sign: a bracketing box $X = [\mathrm{lo}, \mathrm{hi}]$ with
certified opposite signs of $D$ at its ends (certified bisection;
midpoints are outward-rounded balls --- for non-dyadic parameters
the registered exact-dyadic midpoints are not representable --- a
strictly conservative refinement, since wider endpoint balls make
every certificate harder, never easier);
transversality $\inf_{x \in X} D'(x) > 0$ by interval evaluation
over the hull; window inclusion $X \subset [x_\dagger, x_+]$; and
both sides of the scale law.  The ratio is evaluated by the
backward recurrence
$\rho_\alpha = x/(2(\alpha+1)+x\rho_{\alpha+1})$ seeded at order
$\nu+N$ with the \emph{interval} $[0,1]$ --- valid by positivity
and the $c=\tfrac12$ bound of \cite{ErikssonC4}, both outside the
forbidden theorem set --- so no exponentially large value of
$I_\nu$ is ever formed; the seed contraction is certified a
posteriori by the final ball width, and a too-wide ball can only
produce an undecided verdict, never a wrong one.  All $30$ pairs pass all four
certificates at the first ladder rung ($128$ bits); an independent
ordinary high-precision pass (labelled \emph{verified}, not
certified) agrees at all $30$.  Two illustrative rows: at
$(100, 0.501)$ the certified box $[50250.8995\ldots]$ (width
$\sim 3.5\times10^{-14}$) clears the threshold
$x_\dagger = 50250.3995\ldots$ with certified
$D'~\mathrm{ball} = 3.95\times10^{-13} \pm 4.1\times10^{-16} > 0$;
at $(\tfrac12, 0.501)$ the box sits at
$x_* = 501.000000000000000\ldots$, reflecting the closed
(hyperbolic-cotangent) form of $\rho_{1/2}$: by the calibration
\eqref{eq:calibration}, the crossing equation is
$\psi_{1/2}(x_*) = 2(\nu+c)$, and
$\psi_{1/2}(x) = 2 + 1/(x-1)$ up to exponentially small error, so
$x_* \approx 1 + \tfrac{1}{2c-1} = \tfrac{2c}{2c-1} = 501$ --- a
paper-level remark; the certified box needs no such structure.  The companion
is \emph{not} part of the proof, and the proof does not cite it.

\section{Related work}\label{sec:related}

The comparison below follows the development's pre-registered
collation; the antecedents are compared claim by claim.

\textbf{Ruiz-Antol\'in--Segura \cite{RS16}} is the exact antecedent
of the \emph{family} and of \emph{qualitative existence}: under
$\mu = \nu+1$, $\alpha = 2c-1$, their $b_\alpha$ family is
$\Bfam_{\nu,c}$; they identify the sharp ends ($\alpha \le 0$,
i.e.\ $c \le \tfrac12$, at $x \to \infty$; $\alpha \ge 1$, i.e.\
$c \ge 1$, at $x \to 0$) and observe that for intermediate
parameters the graphs must cross (asymptotic end signs).  We
declare expressly: the family and the mere existence of some
crossing are prior art.  \cite{RS16} contains no uniqueness,
orientation, threshold, tangency exclusion, or quantitative
localization for the intermediate members.

\textbf{Hornik--Gr\"un \cite{HG13}} is the closest antecedent in
spirit: it characterizes which two-parameter Amos-type expressions
$G_{\alpha,\beta}$ are upper or lower bounds, and describes minimal
upper bounds as those tangent to the ratio at a single point.  The
statements are about the \emph{boundary of the bounding region} in
the $(\alpha,\beta)$-plane; classifying the fixed diagonal
$G_{\alpha,\alpha}$-type member $\Bfam_{\nu,c}$ --- which for
$c \in (\tfrac12,1)$ is \emph{not} a bound at all --- is a
different question.  Parameter by parameter: \cite{HG13} yields, for
the diagonal, membership/non-membership in the bounding classes; it
does not state, for non-bounding diagonal members, the unique
simple crossing, its transversality, the explicit threshold
$x_\dagger$, bilateral orientation, or a scale law.

\textbf{Segura \cite{Segura23}} studies bounds
$(\alpha+\sqrt{\beta^2+\gamma^2x^2})/x$-type with coefficients
\emph{optimized as functions of a prescribed tangency point}: there
the bounding function is a global bound whose graph is tangent to
the ratio (osculatory contact by construction).  Here the family
member is \emph{fixed in advance}, is not a bound, the meeting is
transversal, and tangential contact is \emph{excluded by theorem}.
Close in machinery, not equivalent in statement.

\textbf{Hornik--Gr\"un \cite{HG24}} proves single-sign-change
results \emph{between two bounds} within their generalized family
(bound-versus-bound comparisons); the present single-crossing
statement is \emph{ratio-versus-family-member}, a different pair of
curves --- we flag the terminological near-collision explicitly.

\textbf{Segura \cite{Segura21}} (Riccati/nullcline monotonicity
methods) and \textbf{Salazar \cite{Salazar26}} (Bernstein-positivity
certificates for Riccati reductions) supply methodology adjacent to
ours; neither presents the quantitative diagonal crossing package.

On the formalization side, the companion papers
\cite{ErikssonC2,ErikssonC3,ErikssonC4} formalized the
$c = \tfrac12$ bound and its consumers; we know of no other
machine-checked development of Amos-type bounds or of any crossing
classification.  The recurrence and derivative identities are
\cite[10.29]{DLMF}.  The development builds on Mathlib
\cite{mathlib20}.

\section{Reproducibility}\label{sec:repro}

Toolchain: Lean \lean{leanprover/lean4:v4.29.0-rc6}; Mathlib pinned
to commit
\begin{center}\small
\lean{07642720480157414db592fa85b626dafb71355b}
\end{center}
\noindent (lakefile and manifest agree); \lean{tectonic} 0.15.0;
\lean{python-flint} 0.9.0 for the certified companion.

\medskip
\noindent Recorded build facts: at the final mathematics commit
\begin{center}\small
\lean{f1c87fef} \emph{(short id; the release manifest records the
full hashes)}
\end{center}
\noindent \lean{lake build AmosClosure} completed successfully
(\textbf{8175 jobs}); the axiom oracle printed, for all
\textbf{110} registered statements (the 74 of the \lean{c4-v1.0}
release registry plus the 36 of this paper), exactly
\begin{center}
\lean{[propext, Classical.choice, Quot.sound]}.
\end{center}
\noindent One proof-engineering disclosure: the division-free
identity inside the tangency-exclusion proof
(\lean{crossing\_no\_critical\_contact}) is closed by
\lean{linear\_combination} with cofactor $K = -2c^2(2\nu+1)$; the
cofactor was found by polynomial division (SymPy \lean{div}, exact,
remainder $0$) against the \lean{ring\_nf}-normalized goal, and is
\emph{replayed entirely inside Lean's ring normalizer} --- the
computer-algebra system plays no trusted role.  The full run
history of this development --- seven green build transcripts and
the failed-attempt record, each committed the day it was produced
with its diagnosis in the commit message; the certified companion
script-and-transcript pair; and eight charter amendments recording
every external audit round, every score checkpoint, and every
method deviation (including the companion's anchor deviation and
depth correction, both registered \emph{before} the certified run)
--- is committed under \lean{scripts/} and
\lean{docs/C5-CHARTER.md}; per the repository's audit protocol it
is archival material and plays no role in the mathematics above.

\medskip
\noindent Headline-theorem-to-artifact map (all in module
\lean{AmosClosure}; every row carries a direct oracle witness; the
74 inherited statements are mapped in
\cite{ErikssonC2,ErikssonC3,ErikssonC4}):

\begin{center}
\footnotesize
\begin{tabular}{@{}llll@{}}
\toprule
Result & Lean name & File & Status \\
\midrule
family, positivity, Eq.~\eqref{eq:calibration} & \lean{amosFamily}, \lean{amosFamily\_calibration} & \lean{AmosFamily.lean} & exact \\
Eq.~\eqref{eq:nullcline} + trichotomy & \lean{riccatiQReal\_amosFamily} etc. & \lean{AmosFamily.lean} & exact \\
sufficiency half & \lean{amosFamily\_upper\_of\_le\_half} & \lean{AmosFamily.lean} & exact \\
ratio recurrence & \lean{besselIReal\_ratio\_recurrence} & \lean{AmosLowerReal.lean} & exact \\
lower bound $L_\nu$ & \lean{besselLowerReal\_holds} & \lean{AmosLowerReal.lean} & exact \\
witness $x_0(c)$ & \lean{amosFamily\_lt\_amosLower\_witness} & \lean{AmosLowerReal.lean} & exact \\
Thm.~\ref{thm:classification} & \lean{amosFamily\_uniform\_upper\_iff} & \lean{AmosLowerReal.lean} & exact \\
Thm.~\ref{thm:lowerhalf} & \lean{amosFamily\_lower\_of\_one\_le} & \lean{AmosCrossing.lean} & exact \\
Thm.~\ref{thm:turan} & \lean{besselIReal\_turan} & \lean{AmosCrossing.lean} & exact \\
window (\S\ref{sec:window}) & \lean{amosFamily\_crossing\_exists} & \lean{AmosCrossing.lean} & exact \\
Eq.~\eqref{eq:slope}, threshold & \lean{crossing\_hasDerivAt}, \lean{crossing\_threshold} & \lean{AmosCrossing.lean} & exact \\
uniqueness, orientation & \lean{crossing\_unique}, \lean{crossing\_orientation\_*} & \lean{AmosCrossing.lean} & exact \\
scale law & \lean{crossing\_scale} & \lean{AmosCrossing.lean} & exact \\
Thm.~\ref{thm:nocontact} & \lean{crossing\_no\_critical\_contact} & \lean{AmosTangency.lean} & exact \\
Thm.~\ref{thm:crossing} package & \lean{amosFamily\_global\_crossing} & \lean{AmosTangency.lean} & exact \\
\S\ref{sec:companion} grid & \lean{c5\_crossing\_arb.py} + transcript & \lean{scripts/} & certified \\
\bottomrule
\end{tabular}
\end{center}

\medskip
\noindent From a clean clone at the release revision (on Windows,
clone with \lean{core.autocrlf=false} --- LF bytes reproduce the
companion script's runtime self-hash --- and
\lean{core.longpaths=true}):
\begin{verbatim}
lake exe cache get
lake build AmosClosure
lake env lean AmosClosure/Oracle.lean
python scripts/c5_crossing_arb.py
tectonic -X compile papers/c5-crossing/c5_crossing.tex
\end{verbatim}
{\sloppy
At the release revision, canonical artifact hashes will be
recorded --- tag-scoped, in \lean{RELEASE-MANIFEST.md} in this
paper's directory --- committed together with the release revision
of this manuscript; compare against
\lean{git show <rev>:path}, never a worktree checkout.  The
annotated tag \lean{c5-v1.0} will mark that release commit after
the pre-release audits close, with the audit registry in
\lean{docs/C5-CHARTER.md}; the release-tree audit is recorded in
the manifest, post-tag by sequencing, with any fixes in post-tag
commits; tags never move.\par}

\section{Conclusion}

For the diagonal Amos-type family at real order, the pinned formal
development now contains the parameter classification ---
uniform upper bound iff $c \le \tfrac12$, uniform lower bound for
$c \ge 1$, with explicit rational witnesses --- and, in the
intermediate regime $\tfrac12 < c < 1$, a global quantitative
crossing theorem: exactly one crossing, transversal, strictly above
an explicit threshold, with bilateral orientation, a constructive
window, a two-sided scale law, and degenerate contact excluded by
an exact curvature identity.  Limitations: the in-core series
definition only, and $\nu \ge 0$ (Section~\ref{sec:scope}); the
family, the classification as mathematics, and the crossing
phenomenon are prior art \cite{RS16} --- the quantitative crossing
package is the candidate-new content, stated with the priority
prudence of Section~\ref{sec:scope}.  Natural next steps: the
analogous classification for the second-kind ratio
$K_{\nu+1}/K_\nu$ (where \cite{Segura23} suggests the tangency
structure differs), asymptotics of $x_*(\nu,c)$ beyond the scale
law, and upstreaming the modified-Bessel API.

\begin{thebibliography}{99}

\bibitem{Amos74} D.~E.~Amos,
\emph{Computation of modified Bessel functions and their ratios},
Math. Comp. \textbf{28} (1974), 239--251.
doi:10.1090/S0025-5718-1974-0333287-7.

\bibitem{Segura11} J.~Segura,
\emph{Bounds for ratios of modified Bessel functions and associated
Tur\'an-type inequalities},
J. Math. Anal. Appl. \textbf{374} (2011), 516--528.
doi:10.1016/j.jmaa.2010.09.030.

\bibitem{HG13} K.~Hornik and B.~Gr\"un,
\emph{Amos-type bounds for modified Bessel function ratios},
J. Math. Anal. Appl. \textbf{408} (2013), 91--101.
doi:10.1016/j.jmaa.2013.05.070.

\bibitem{RS16} D.~Ruiz-Antol\'in and J.~Segura,
\emph{A new type of sharp bounds for ratios of modified Bessel
functions}, J. Math. Anal. Appl. \textbf{443} (2016), 1232--1246.
doi:10.1016/j.jmaa.2016.06.011.  arXiv:1606.02008.

\bibitem{Segura21} J.~Segura,
\emph{Monotonicity properties for ratios and products of modified
Bessel functions and sharp trigonometric bounds},
Results Math. \textbf{76} (2021), no.~4, Paper No.~221.
doi:10.1007/s00025-021-01531-1.  arXiv:2105.02524.

\bibitem{Segura23} J.~Segura,
\emph{Simple bounds with best possible accuracy for ratios of
modified Bessel functions},
J. Math. Anal. Appl. \textbf{526} (2023), no.~1, 127211.
doi:10.1016/j.jmaa.2023.127211.  arXiv:2207.02713.

\bibitem{HG24} K.~Hornik and B.~Gr\"un,
\emph{Generalized Amos-type bounds for modified Bessel function
ratios}, Math. Inequal. Appl. \textbf{27} (2024), 775--787.
doi:10.7153/mia-2024-27-55.

\bibitem{Salazar26} D.~S.~P.~Salazar,
\emph{Riccati reductions for modified Bessel ratios: Bernstein
positivity, exact certificates, and transfer obstructions},
arXiv:2607.05538, 2026.

\bibitem{BP13} \'A.~Baricz and S.~Ponnusamy,
\emph{On Tur\'an type inequalities for modified Bessel functions},
Proc. Amer. Math. Soc. \textbf{141} (2013), 523--532.
doi:10.1090/S0002-9939-2012-11325-5.  arXiv:1010.3346.

\bibitem{DLMF} NIST Digital Library of Mathematical Functions,
F.~W.~J.~Olver et al., eds., \url{https://dlmf.nist.gov/}.

\bibitem{ErikssonC2} L.~Eriksson,
\emph{One Amos bound, three consumer sites: machine-checked
Bessel-ratio calculus for lattice gauge expansions}, companion
paper, same repository, release \lean{c2-v1.2.1};
ai.viXra.org:2607.0033, July 2026.

\bibitem{ErikssonC3} L.~Eriksson,
\emph{A machine-checked proof of the Amos bound for modified Bessel
function ratios}, companion paper, same repository, release
\lean{c3-v1.0.1}; ai.viXra.org:2607.0032, July 2026.

\bibitem{ErikssonC4} L.~Eriksson,
\emph{A machine-checked proof of an Amos-type bound for modified
Bessel ratios at real order}, companion paper, same repository,
release \lean{c4-v1.0}; ai.viXra.org:2607.0030, July 2026.

\bibitem{mathlib20} The mathlib Community,
\emph{The Lean mathematical library}, in: Proc. CPP 2020, ACM, 2020,
pp. 367--381. doi:10.1145/3372885.3373824.

\end{thebibliography}

\end{document}
